\documentclass[11pt]{article}

% This file will be kept up-to-date at the following GitHub repository:
%
% https://github.com/automl-conf/LatexTemplate
%
% Please file any issues/bug reports, etc. you may have at:
%
% https://github.com/automl-conf/LatexTemplate/issues

\usepackage{microtype} % microtypography
\usepackage{booktabs}  % tables
\usepackage{url}  % urls
\usepackage{csquotes}

% AMS math
\usepackage{amsmath}
\usepackage{amsthm}

\newcommand{\netpool}{\vert \mathcal N \vert}

% With no package options, the submission will be anonymized, the supplemental
% material will be visible, and line numbers will be added to the manuscript.
%
% To compile a non-anonymized camera-ready version, add the [final] option:
%
% \usepackage[final]{automl}
%
% To compile for the non-archival content track, add the [shortpaper] option:
%
% \usepackage[shortpaper]{automl}
% \usepackage[final,shortpaper]{automl}
%
% To compile a single-blind submission, e.g. for the ABCD track, add the
% [revealauthors] option:
%
% \usepackage[revealauthors]{automl}
%
% To compile a non-anonymized pre-print version that can be uploaded to arXiv,
% add the [preprint] option:
%
% \usepackage[preprint]{automl}
%
% The [hidesupplement] option may be used to hide the supplementary material from any of the
% above commands, e.g.:
%
% \usepackage[hidesupplement,final,shortpaper]{automl}

\usepackage{automl}

% You may use any reference style as long as you are consistent throughout the
% document. As a default, we suggest author--year citations; for bibtex and
% natbib you may use:

\usepackage{natbib}
\bibliographystyle{apalike}

\newcommand{\submissiontitle}{Learning to Search: HW-NAS Probing Few Networks}

\title{\submissiontitle}
\author[1]{\nameemail{Francesco Capuano}{francesco.capuano@ens-paris-saclay.fr}}

% The syntax for adding an author is
%
% \author[i]{\nameemail{author name}{author email}}

% where i is an affiliation counter. Authors may have
% multiple affiliations; e.g.:
%
% \author[1,2]{\nameemail{Anonymous}{anonymous@example.com}}
%
% the list might continue:
% \author[2,3]{\nameemail{Author 2}{email2@example.com}}
% \author[3]{\nameemail{Author 3}{email3@example.com}}
%
% if you need to force a linebreak in the author list, prepend an \author entry
% with \\:
%
% \author[1,2,3]{\\\nameemail{Author 4}{email4@example.com}}
%
% Specify corresponding affiliations after authors, referring to counter used in
% \author:
%
\affil[1]{École Normale Supérieure Paris-Saclay}
%
% the list might continue:
% \affil[2]{Institution 2}
% \affil[3]{Institution 3}
%
% this system may also be used to add additional footnotes, e.g. to indicate
% equal contributions. Here is an example:
%
% \author[1,$\ast$]{\nameemail{Author 1}{email1@example.com}}
% \author[1,$\ast$]{\nameemail{Author 2}{email2@example.com}}
% \author[2]{\nameemail{Author 3}{email3@example.com}}
%

% define PDF metadata, aids in accessibility of the resulting PDF
\hypersetup{%
  pdfauthor={Francesco Capuano}, % will be reset to "Anonymous" unless the "final" package option is given
  pdftitle={\submissiontitle},
  pdfsubject={MetaRL for NAS},
  pdfkeywords={AutoML, Reinforcement Learning, NAS}
}

\begin{document}

\maketitle

\begin{abstract}
    Existing hardware-aware NAS (HW-NAS) methods typically assume access to precise information about the target device, either via analytical approximations of post-compilation latency or through learned latency predictors. Such approaches limit generalization and can introduce estimation errors that are detrimental in risk-sensitive applications.

We propose a two-stage HW-NAS framework, in which we first \emph{meta-learn} an architecture controller on a distribution of synthetic devices. At test time, the trained controller deploys directly to the target device without relying on any pre-collected information, and rather only exploits direct interactions. In particular, the controller designs an architecture for the target device by interacting with it through a small number of high-fidelity latency measurements, thanks to the meta-learning phase. Thus, our approach dispenses with amortized latency estimation at deployment, reducing model-error risks and enabling robust HW-NAS in scenarios where one cannot trust approximate latency sources.

To control computational costs, we train our controller with training-free accuracy proxies, allowing us to scale the meta-training phase without incurring the overhead of full network training. We benchmark on HW-NATS-Bench, demonstrating that our method \emph{generalizes to unseen devices} and searches for latency-efficient architectures by in-context adaptation using only a few real-world latency evaluations at test-time.
\end{abstract}

\section{Introduction}
Neural Architecture Search (NAS) has emerged as a powerful paradigm for automating neural network design. While significant progress has been made in improving sample efficiency and final performance, most NAS methods remain specialized to a single deployment target. This limits their practicality in real-world scenarios where architectures must be deployed across a diverse set of hardware platforms, each with distinct computational constraints and performance tradeoffs.

To address this limitation, we focus on the problem of \textit{multi-device NAS}, where the goal is to discover architectures that are optimized for hardware-specific metrics (e.g., latency, memory footprint) and that generalize to unseen devices with minimal adaptation cost. Current hardware-aware NAS approaches \cite{wu2019fbnet, cai2019proxylessnas} rely on performance measurements collected from a fixed device and fail to generalize when the deployment conditions change. In contrast, we propose to train a NAS controller that is exposed during training to a diverse distribution of hardware devices. This enables the controller to learn generalizable priors over device-conditioned architecture search trajectories.

Our key insight is to cast hardware-aware NAS as a meta-learning problem. We meta-train a reinforcement learning (RL) controller across multiple devices, each sampled with its own latency and task constraints. The controller learns to maximize a training-free proxy for downstream accuracy, enabling efficient evaluation without full model training. At test time, the controller adapts to new hardware targets using only a few queries, achieving strong performance on devices never seen during meta-training.

This work makes the following contributions:
\begin{itemize}
    \item We formalize the task of multi-device neural architecture search and propose a meta-learning framework to address it.
    \item We design a domain-randomized RL training procedure that enables the NAS controller to generalize across heterogeneous devices.
    \item We empirically validate our method on multiple edge devices, demonstrating its ability to perform few-shot NAS on previously unseen hardware platforms.
\end{itemize}

By enabling cross-device generalization, our approach offers a scalable path toward deployable NAS in practical settings, where architecture specialization must be fast, accurate, and hardware-aware.

\section{Background}
\label{background}
\subsection{Hardware-Aware Neural Architecture Search (HW-NAS)}

\paragraph{Neural Architecture Search (NAS)}
Designing Deep Neural Networks (DNNs) by hand is cumbersome and brittle process, originally carried out by highly-specialized human experts, and later automatized via Neural Architecture Search (NAS).
Instead of relying on the combination of human ingenuity, heuristics, and trial-and-error to design DNNs, NAS provides an automated and systematic strategy to discover architectures optimizing a specified objective. 
In its seminal work,~\citet{zoph2016neural} demonstrated the feasibility of designing DNNs automatically via Reinforcement Learning (RL), using post-training performance metrics as reward signal.
While effective,~\citep{zoph2016neural}'s pioneering work most notably demanded thousands of GPU hours---a computational demand amply exceeding the capabilities of most institutions.
Subsequent research introduced more computationally efficient approaches, including weight-sharing paradigms~\citep{brock2017smash} and gradient-based methods~\citep{liu2018darts}.
A major goal of NAS is to reduce human intervention in selecting specialized architectures for novel tasks, data modalities, or deployment conditions. By systematically exploring a well-defined search space, NAS not only identifies high-accuracy networks but can also reveal unusual architectures that depart from human-engineered patterns. 
In the years following~\citet{zoph2016neural}'s work, NAS has expanded to address a variety of tasks within Deep Learning (DL), including language modeling, and sequence-to-sequence learning~\citep{elsken2019neural}. 
Despite these advances, the overall \emph{cost of search} remains a key concern, especially when architectures must be verified experimentally for downstream performance. In practice, it is common to search for architectures within cell-based searchspaces~\citep{zoph2018learning}, hereby reducing DNN design to identifying an optimal cell within a possibly very large searchspace of all possible combinations.
Under reasonable assumptions, the computational burden of NAS \(C \) can be modeled as the product of (i) the number of candidate architectures probed \( \#\text{networks}\), and (ii) the average \emph{cost to evaluate one network}, \( c \), yielding \( C = \sum_i^{\mathcal N} c_i \approx \netpool c \), for a pool of probed networks, \(\mathcal N \). 
The cost of many early methods~\citep{zoph2016neural,baker2016designing,real2017large} was mainly dominated by \( c \), as they required \emph{training} candidate architectures on a target task to measure performance. Several strategies have been introduced to make evaluation more tractable. 
Training a single hypernetwork, for instance, amortizes \(C \) by training a single overparameterized model from which many sub-architectures can be sampled~\citep{brock2017smash, pham2018efficient,cai2019once}, hereby reducing \(C \) via reducing \( c \). However, training such a hypernetwork is still prohibitively expensive in scenarios where computational resources are not abundant, and does not have access to thousands of GPU hours~\citep{cai2019once} when searching large spaces of candidate architectures.

\paragraph{Training-free (TF) metrics}
To tackle these limitations, the research community developed \emph{training-free metrics} (TF metrics)~\citep{mellor2021neural,jacot2018neural,chen2021neural,chen2023understanding}, which allow to drastically reduce the and cost and resources needed to evaluate candidate architectures by scoring them at initialization, without training.
Thus, TF metrics serve as proxies for post-training performance, guiding search toward promising architectures without requiring any update of any architecture parameters. In particular, in this work we focus on~\emph{NASWOT}~\citep{mellor2021neural}, which quantifies the expressivity of an untrained network by examining its weight initialization and the distinct activation patterns it can generate. By analyzing how the network divides input data---by measuring the number of "linear regions" formed by decision boundaries at initialization---NASWOT produces a one-shot score correlating with final performance.
To search large spaces of candidate architectures, \citet{cavagnero2023freerea} propose complemeting NASWOT with \emph{LogSynflow}~\citep{cavagnero2023freerea, tanaka2020pruning}, correlating the performance of a candidate design with the gradient flow of a loss function specifically designed to reflect the model's capacity.
Searching architectures using TF metrics substantially reduces the cost of probing each candidate architecture \( c \). In turn, TF metrics prove particularly useful to guide possibly sample-inefficient algorithms, such as RL~\citep{zoph2016neural}, allowing RL to rival more sample-efficient methods like evolutionary algorithms~\citep{real2017large} for cell-based NAS.
Cell-based NAS algorithms are usually compared on cell-based search spaces such as NAS-Bench-101~\citep{ying2019bench} or NATS-Bench~\citep{dong2021nats}.

\paragraph{Hardware-Aware NAS (HW-NAS)}
While traditional NAS typically focuses on metrics such as validation accuracy or model size, hardware-aware NAS (HW-NAS) seeks to incorporate hardware performance indicators characteristic of the deployment scenario of the model while designing DNNs~\citep{benmeziane2023efficient}. 
Work in HW-NAS typically considers optimizing post-compilation latency~\citep{wu2019fbnet}, memory footprint~\citep{king2025micronas} or energy consumption~\citep{speckhard2022neural}. In this work we exclusively focus on \emph{post-compilation latency}, as it is crucial when deploying DNNs to edge devices and resource-constrained environments where low latency inference is paramount, such as autonomous robotics and internet of things (IoT) devices. 
Taking into account hardware performance when designing DNNs introduces another dimension of complexity, as architectures that perform well on a given downstream task on the high-end devices used to train them---typically, high-end GPUs or even high-performance computing (HPC) infrastructures---might be infeasible due to unacceptable latency or resource overhead in different deployment scenarios. 
Nonetheless, designing architectures for real-world use does require taking into account the specifities of the deployment scenario, such as the device used, or the expected workload. HW-NAS seeks to address precisely this very challenge, complementing NAS' automated approach to DNN design with considerations for hardware constraints and performance metrics specific to the deployment environment~\citep{li2021hw}.
However, designing a comprehensive HW-NAS strategy concerned with downstream real-world performance of trained DNNs is a challenging endeavour. 
First, measuring hardware performance at inference-time typically involves low-level compiling and deploying candidate devices, which can prove challenging and for non-domain experts, other than being costly~\citep{li2021hw}.
Second, hardware performance profiles can vary significantly across devices due to diverse hardware pipelines, compilers, and operator libraries~\citep{li2021hw,benmeziane2023efficient}.
In practice, the same high-level logical operations---e.g., a 3x3 convolutional filter---are compiled very differently on different hardware platforms~\citep{li2021hw,benmeziane2023efficient,laube2022expect}. In turn, this makes it nontrivial to identify good architectures across diverse devices.
HW-NAS-Bench~\citep{li2021hw} was proposed to promote reproducibility and expedite research in hardware-aware design for multi-device setting. 
Similarly to NATS-Bench~\citep{dong2021nats}, HW-NAS-Bench provides a curated set of architectures along with measured latency and energy consumption across multiple real-world hardware targets, ranging from commercial edge devices (Raspberry Pi, Google Pixel3) to specialized accelerators (FPGA, Eyeriss).

\subparagraph{Measuring latency for HW-NAS}
Measuring actual inference latency often proves crucial in hardware-aware NAS, yet directly compiling and timing each candidate architecture can be a nontrivial and costly process~\citep{li2021hw}. In principle, obtaining a ground-truth latency measurement for a specific model-device pair entails performing low-level optimizations. Differences in compiler toolchains or library support exacerbate this burden, resulting in latency measurements that can vary across hardware generations or even across software updates.
As an alternative, many HW-NAS approaches rely on \emph{approximations} of post-compilation latency. 
A popular choice is using lookup tables (LUT), modelling architectures \( a = \{o_1, o_2, \dots \} \) as a set of fundamental operations, and aggregating the per-operation runtimes \( t_d(o_i), \ o \in \mathcal O \) to obtain the network's latency \( \ell_d(a) \) for a given device \( d \). 
Typically, operation-wise runtimes are typically aggregating as \(\ell (a) \approx \hat{\ell}(a) = \sum_{o \in a} t_d(o) \)~\citep{wu2019fbnet,cai2018proxylessnas}. 
LUTs are fast to use at query time, and cheap to obtain even for very large search spaces. 
For instance, while FBNet~\citep{wu2019fbnet} contains \(\sim10^{21} \) architectures, obtaining a LUT for FBNet only requires deploying its 9 fundamental operations.
However, LUT deliberately fail to capture interactions among operations within a given network, which results in poor correlation between approximated and true latency values~\citep{laube2022expect}.
A more flexible strategy is to learn a \emph{predictor model} for latency, where an architecture is encoded using a feature vector \( h(a) \), and the associated latency \( \ell(a) \) is regressed using supervised learning on pairs of \( \tau = \{h(a_i), \ell(a_i)\}_{i=1}^{K} \)~\citep{laube2022expect}. Once trained, the predictor can quickly infer latency for unseen architectures, thereby accelerating the NAS loop.
\citet{laube2022expect} report the Kendall-Tau correlation (KT) between true and approximated latency values of various predictors versus LUT varying the number of examples needed to train a regression model \( K \), and report any of the predictors used---ranging from gradient-boosted trees to MLPs---drastically outperforms LUT's KT correlation of 0.6 with as little as \(K=60\) examples, reaching KT values of 0.8 with \( K = 500 \). 
In turn, \citet{laube2022expect} show that \emph{(ii)} LUT introduce severe misestimations of true latency, with unsatisfactory correlation with real-world latency and \emph{(ii)} predictors trained on a few hundreads of examples outperform them significantly.
However, such predictive methods introduce approximation errors making them unreliable in risk-averse deployments such as autonomous robotics, aerospace, or industrial monitoring.
Further, both LUTs and predictors assume substantial upfront data collection---building the table or training the model requires a nontrivial set of measured (architecture, latency) samples. When the target hardware is scarce or each compilation is expensive, these preliminary steps may already pose a major overhead. Furthermore, if practitioners suspect that distribution shifts might occur (e.g., minor firmware updates), a learned or tabulated model may become out-of-date, again risking suboptimal selections.

In contrast, our approach circumvents the use of any learned predictor or LUT by constraining the total number of architectures probed to a small but carefully selected set, hereby limiting \( \netpool \) at test time. 
Because we only need to compile and measure a handful of candidate networks, we can reliably obtain high-fidelity latency evaluations for each one, for small \( C \). 
This avoids introducing approximation errors into the search loop and reduces the risk of committing to an inaccurately estimated architecture in mission-critical deployments.

\subsection{Reinforcement Learning (RL)}
Reinforcement learning (RL) is an area of machine learning concerned with training agents to act optimally in an environment that provides feedback via a reward signal~\citep{sutton1998reinforcement}. 
In general, RL agent act on environment's states \(s_t \in \mathcal S \), to taking an action \(a_t \in \mathcal A \). 
For each interaction \(t=1, \dots, T \) the agent receives a reward \(r_t \in \mathbb R\) reflecting the desirability of taking action \( a_t \) when faced with \( s_t \), in light of the agent's ultimate goal of maximizing the total sum of (discounted) rewards \(\sum_t^T \gamma^t r_t \).
Over time, RL agents attempt to learn a conditional probability distribution \(\pi: \mathcal S \times \mathcal A \mapsto [0, 1] \) ("policy", \(\pi(a \vert s)\)) following this objective of cumulative reward maximization.
Over discrete state and action spaces---typical in cell-based NAS where any \( s \in \mathcal S \) represents a given cell and actions \( a \in \mathcal A \) discrete modifications of candidate architectures---popular RL algorithms include trust-region approaches, such as Trust Region Policy Optimization (TRPO)~\citep{schulman2015trust} and Proximal Policy Optimization (PPO)~\citep{schulman2017proximal}.
Unlike REINFORCE~\citep{zoph2016neural}, trust-region approaches such as PPO iteratively learn \( \pi \) bounding the relative difference between successive polices. In this, trust-region approaches aim at avoiding too drastic approaches from which on-policy RL risks to not recover otherwise.

While RL-based NAS set an important precedent, it was quickly recognized that repeatedly training every candidate network rendered the search protocol too costly. Thus, researchers moved towards alternative solutions that attempt to reduce computational demand by limiting the number of networks probed \( \netpool \), such as evolutionary algorithms~\citep{real2017large,cavagnero2023freerea}.
While evolutionary algorithms (EA) have shown promise for single-device NAS~\citep{real2017large,cavagnero2023freerea}, they remain ill-suited to scenarios in which multiple target devices must be supported, especially when each device differs substantially in its architectural or compiler-level constraints. 
Standard EA-based NAS pipelines re-initialize or re-search from scratch whenever the hardware setting changes, thereby incurring overhead. Even when combined with hardware metrics, EAs is inherently ill-posed to transfer previously acquired knowledge to a new device without discarding substantial parts of the discovered solutions.

In contrast, RL has exhibited strong cross-domain adaptability in fields where cross domain adaptability is key, such as robotics~\citep{kober2013reinforcement,akkaya2019solving}, and we thus argue RL's capacity for \emph{transfer learning} is especially attractive in HW-NAS. 
Rather than fully re-searching architectures for each new deployment device, one can meta-train an RL controller on a distribution of possible hardware profiles and re-use this controller whenever a novel device or updated firmware is encountered. 
Crucially, this approach avoids the need for separate predictor models per device; the policy itself encodes adaptation strategies, sampling only a few real-world measurements to confirm the final architecture design.
In this, RL's cost is limited not by limiting \( c \), but rather \( \netpool \) at test-time.

\subsection{Related Works}

\paragraph{Once-for-All (OFA)}
A pioneering approach in low-cost HW-NAS Once-for-All (OFA)~\citep{cai2019once}. OFA amortizes NAS' costs \( C \) by training a single over-parameterized hypernetwork which is pruned via EA~\citep{real2017large} into "child" sub-networks, with each subnetwork inheriting weights from the super-network without requiring any additional training epochs~\citep{brock2017smash}.
In principle, this allows cheap evaluation of the candidate architectures at test-time (\( c \downarrow \implies C \downarrow \)), by amortizing the cost of training the hypernetwork---1.2k GPU hours on NVIDIA V100 GPUs---on multiple searches. OFA also allows for HW-NAS, using LUT to guide EAs in selecting architectures fitting specific computational constraints. 
However, OFA's success relies on the costly super-network pre-training phase, which requires access to GPU hours simply unavailable in most scenarios. Further, sub-network selection process is also subject to approximation errors in hardware metrics, consequent to using LUT to estimate latency, which can yield sub-optimal sub-network choices~\citep{laube2022expect}.

\paragraph{Hardware-adaptive Efficient Latent Predictor (HELP)}
Hardware-adaptive Efficient Latent Predictor (HELP)~\citep{lee2021hardware} proposes a similar approach of keeping \( C \) manageable, by learning a meta-predictor to estimate post-compilation latency for unseen architectures, thus focusing on \( c \downarrow \implies C \downarrow \) too. Rather than measuring the latency of every candidate architecture on the physical device, HELP trains a surrogate model to map architecture encodings to hardware costs. Then, by assessing as little as 10 architectures on the target device, HELP adapts the meta-learned predictor, thus severily limiting the \emph{per-probe cost} \(c\) of each architecture. 
Yet, constructing HELP's predictor entails collecting a non-trivial set of real-world measured samples \(\tau_d=\{h(a_i), \ell(a_i)\}_{i=1}^K\) across possibly many different devices, which can prove challenging for non-experts in hardware~\citep{li2021hw}. Further, if the deployment scenario---the hardware platform or software environment---changes, the predictor may require retraining or fine-tuning. Lastly, in risk-sensitive deployments even modest misestimation of latency can invalidate guarantees on latency bounds, hindering applicability of predictor-based approaches such as HELP. 

Both OFA and HELP aim to mitigate \( C \) by limiting the overhead associated with \(c \). OFA amortizes \( c \) by reusing super-network weights and LUT---at the expense of accuracy in latency estimation---while HELP uses a meta-learned latency regressor---at the expense of safety in latency estimation. 
By contrast, our method explicitly mitigates \( C \) by limiting the \emph{number} of high-fidelity probes \(\netpool \) at test-time, sidestepping super-network pre-training or predictive models altogether. 
This design choice inherently avoids approximation errors in hardware metrics, as each of the few final architectures is tested with real measurements. Hence, we shift complexity into a meta-trained RL controller (Section~\ref{sec:method}) that can generalize over device variations learning to design well-performing DNNs across various target platforms.

\section{Method}
\label{sec:method}

% High-level overview:
% 1. Problem Setting / Notation
% 2. RL-based approach for multi-device adaptation
% 3. Architecture encoding & search space
% 4. Training-free metric integration
% 5. Meta-training procedure & domain randomization
% 6. Deployment stage: real-world latency measurement

In this section, we present the details of our hardware-aware NAS framework that meta-learns a reinforcement learning (RL) controller to operate across multiple devices. We first formalize the problem setting and notation (\S\ref{sub:problem_setup}), then describe the proposed RL-based search procedure (\S\ref{sub:rl_search}). Next, we illustrate how we encode architectures and integrate training-free metrics (\S\ref{sub:arch_encoding}), before clarifying our meta-training algorithm and its domain randomization strategy (\S\ref{sub:meta_training}). Finally, we discuss how the controller is deployed and how real-world latency is measured with minimal overhead (\S\ref{sub:deployment_stage}).

\subsection{Problem Setup and Notation}
\label{sub:problem_setup}
We consider a multi-device HW-NAS problem in which the goal is to discover an architecture $a \in \mathcal{A}$ that balances downstream accuracy (or a training-free proxy) and real-world latency on a target hardware device. 
Formally:
\[
    \underset{a \in \mathcal{A}}{\text{maximize}} \quad \mathrm{AccProxy}(a) \quad \text{subject to} \quad \mathrm{Latency}_{\mathrm{real}}(a,d) \leq \tau,
\]
where $\mathrm{Latency}_{\mathrm{real}}(a,d)$ is the measured post-compilation latency of architecture $a$ on device $d$, and $\tau$ is an admissible threshold for that device. 
Our method extends this objective across multiple possible devices $\{d_1, d_2, \dots\}$, each with a different $\tau_i$, allowing us to handle device variation in a single framework.

\subsection{RL-based Search Procedure}
\label{sub:rl_search}
\paragraph{RL Environment Design.}
We frame the architecture selection problem as a sequential decision process, where each state $s_t$ represents a partially built candidate network (or an existing architecture), and actions $a_t$ correspond to modifying or finalizing the architecture. The environment provides a scalar reward, computed from training-free proxies and hardware feedback, guiding the RL policy $\pi_\theta$ to select promising architectures while respecting latency constraints. 

\paragraph{Policy Parameterization.}
The controller $\pi_\theta$ is typically represented by a neural network (e.g., a multi-layer perceptron or LSTM) that outputs a distribution over possible actions given the current architecture encoding $h(s_t)$. We adopt a trust-region or policy-gradient algorithm (e.g., PPO) to iteratively optimize $\theta$. 

\paragraph{Reward Signal.}
The controller’s reward at each timestep aggregates:
\begin{itemize}
    \item \textbf{Proxy Score} (e.g., NASWOT or LogSynflow), quantifying the architecture’s potential accuracy at initialization.
    \item \textbf{Latency Penalty}, based on approximate or synthetic device performance metrics, ensuring that the controller learns to prefer architectures predicted to fit the latency constraint $\tau$.
\end{itemize}
When we deploy on a real device, we refine or finalize the policy’s choice with direct measurements on a \emph{small subset} of candidate architectures.

\subsection{Architecture Encoding and Training-free Metrics}
\label{sub:arch_encoding}
\paragraph{Search Space Representation.}
We adopt a cell-based search space (e.g., NAS-Bench-201) or a more flexible macro-level design, where each layer is selected from a set of possible operations. Each architecture $a \in \mathcal{A}$ can be encoded as a sequence of discrete tokens $\{o_1, o_2, \dots \}$ indicating convolution types, kernel sizes, skip connections, or other hyperparameters.

\paragraph{Training-free Proxy Integration.}
At every RL timestep, the environment computes a training-free metric (NASWOT, LogSynflow, \dots) for the candidate $a$. This quick evaluation provides a gradient-free proxy for final accuracy. Notably, we do \emph{not} train or fine-tune $a$; we only measure the untrained model’s properties. This maintains the low per-probe cost $c$ crucial to scaling RL-based NAS.

\subsection{Meta-Training with Domain Randomization}
\label{sub:meta_training}
\paragraph{Synthetic Hardware Profiles.}
To enable cross-device generalization, we create multiple synthetic device profiles that approximate or bracket the latency distributions of real hardware. Each profile randomly varies operator latencies, memory limits, or pipeline stages, exposing the RL controller to a broad distribution of pseudo-latency signals during training.

\paragraph{Controller Optimization.}
We run RL (e.g., PPO) for multiple episodes on the domain-randomized environment, sampling a different synthetic profile each episode or mini-batch. Over time, $\pi_\theta$ learns robust strategies that do not overfit to a single device. Crucially, this meta-training stage requires \emph{no} real-world device measurements, so we avoid the large overhead of repeated compilation or inference runs.

\paragraph{Few-Shot Adaptation.}
At deployment time, we can optionally allow a small number of real measurements from the target device to fine-tune or guide the controller’s final architecture choices. This step ensures reliability for risk-averse domains by eliminating the approximation gap from synthetic distributions.

\subsection{Deployment Stage: Real-World Latency with Limited Probes}
\label{sub:deployment_stage}
Finally, we describe how the meta-trained controller is used in real-world HW-NAS:
\begin{enumerate}
    \item \textbf{Controller Proposes Candidates.} The RL policy $\pi_\theta$ (trained on synthetic profiles) proposes a small set of architectures $\{a_1, \dots, a_k\}$.
    \item \textbf{Measure True Latency.} Each proposed architecture $a_i$ is compiled and measured on the actual device $d$. Only $k \ll \lvert \mathcal{A}\rvert$ measurements are performed, making the total cost $C$ remain feasible.
    \item \textbf{Select the Final Architecture.} The best architecture that satisfies latency constraints while maximizing the proxy-based performance is chosen. No predictor or LUT is needed, as each candidate’s cost is known exactly.
\end{enumerate}
By limiting $k$, we avoid the overhead of measuring all possible architectures. The policy’s meta-learned knowledge ensures that most proposed designs are high-quality, thus preserving search efficiency. 



\section{Experiments}
\input{sections/04_experiments}

\section{Conclusions}
\input{sections/05_conclusions}


\bibliography{references}

% \newpage
% \section{Submission Checklist}

% All submissions are welcome to include the AutoML submission checklist provided
% below. \textbf{While not required, including the checklist may be beneficial
%   during the reproducibility review process and for maximizing long-term
%   impact.} If you choose to include the checklist, it will not count towards the
% page limit. The submission checklist draws upon related submission checklists:
% %
% \href{https://neurips.cc/Conferences/2022/PaperInformation/PaperChecklist}
%      {the NeurIPS '22 checklist}
% %
% and the
% \href{https://www.automl.org/wp-content/uploads/NAS/NAS_checklist.pdf}
%      {\textsc{nas} checklist}.
% %

% For each question, change the default \verb|\answerTODO{}| (typeset \answerTODO)
% to
% \verb|\answerYes{[justification]}| (typeset \answerYes),
% \verb|\answerNo{[justification]}| (typeset \answerNo), or
% \verb|\answerNA{[justification]}| (typeset \answerNA).
% We recommend including a brief justification to your answer, either by
% referencing the appropriate section of your paper or providing a brief inline
% description.  For example:
% \begin{itemize}
% \item Did you include the license of the code and datasets? \answerYes{See
%   Section~\ref{sec:code}.}
% \item Did you include all the code for running experiments? \answerNo{We include
%   the code we wrote for conducting the experiments, but complete replication
%   depends on proprietary libraries for executing on a private compute
%   cluster. The code therefore is not runnable without modification. To
%   compensate, we provide a runnable but non-parallelized version of the code
%   that could replicate the results at the expense of a greater wall-clock time.}
% \item Did you include the license of the datasets? \answerNA{Our experiments
%   were conducted on publicly available datasets and we have not introduced new
%   datasets.}
% \end{itemize}
% Please note that if you answer a question with \verb|\answerNo{}|, we recommend
% providing an explanation and/or compensation for the omission. For example, if
% you cannot provide complete evaluation code for some reason, you might instead
% provide code for a minimal reproduction of the main insights of your paper.

% Please do not modify the questions and only use the provided macros for your
% answers. Note that the submission checklist does not count towards the page
% limit. If you choose to modify \texttt{instructions.tex}, please delete these
% instructions and only keep the Submission Checklist section heading above along
% with the questions/answers below.

% \begin{enumerate}
% \item For all authors\dots
%   %
%   \begin{enumerate}
%   \item Do the main claims made in the abstract and introduction accurately
%     reflect the paper's contributions and scope?
%     %
%     \answerTODO{}
%     %
%   \item Did you describe the limitations of your work?
%     %
%     \answerTODO{}
%     %
%   \item Did you discuss any potential negative societal impacts of your work?
%     %
%     \answerTODO{}
%     %
%   \item Did you read the ethics review guidelines and ensure that your paper
%     conforms to them? (see \url{https://2022.automl.cc/ethics-accessibility/})
%     %
%     \answerTODO{}
%     %
%   \end{enumerate}
%   %
% \item If you ran experiments\dots
%   %
%   \begin{enumerate}
%   \item Did you use the same evaluation protocol for all methods being compared (e.g.,
%     same benchmarks, data (sub)sets, available resources, etc.)?
%     %
%     \answerTODO{}
%     %
%   \item Did you specify all the necessary details of your evaluation (e.g., data splits,
%     pre-processing, search spaces, hyperparameter tuning details and results, etc.)?
%     %
%     \answerTODO{}
%     %
%   \item Did you repeat your experiments (e.g., across multiple random seeds or
%     splits) to account for the impact of randomness in your methods or data?
%     %
%     \answerTODO{}
%     %
%   \item Did you report the uncertainty of your results (e.g., the standard error
%     across random seeds or splits)?
%     %
%     \answerTODO{}
%     %
%   \item Did you report the statistical significance of your results?
%     %
%     \answerTODO{}
%     %
%   \item Did you use enough repetitions, datasets, and/or benchmarks to support
%     your claims?
%     %
%     \answerTODO{}
%     %
%   \item Did you compare performance over time and describe how you selected the
%     maximum runtime?
%     %
%     \answerTODO{}
%     %
%   \item Did you include the total amount of compute and the type of resources
%     used (e.g., type of \textsc{gpu}s, internal cluster, or cloud provider)?
%     %
%     \answerTODO{}
%     %
%   \item Did you run ablation studies to assess the impact of different
%     components of your approach?
%     %
%     \answerTODO{}
%     %
%   \end{enumerate}
%   %
% \item With respect to the code used to obtain your results\dots
%   %
%   \begin{enumerate}
% \item Did you include the code, data, and instructions needed to reproduce the
%     main experimental results, including all dependencies (e.g.,
%     \texttt{requirements.txt} with explicit versions), random seeds, an instructive
%     \texttt{README} with installation instructions, and execution commands
%     (either in the supplemental material or as a \textsc{url})?
%     %
%     \answerTODO{}
%     %
%   \item Did you include a minimal example to replicate results on a small subset
%     of the experiments or on toy data?
%     %
%     \answerTODO{}
%     %
%   \item Did you ensure sufficient code quality and documentation so that someone else
%     can execute and understand your code?
%     %
%     \answerTODO{}
%     %
%   \item Did you include the raw results of running your experiments with the given
%     code, data, and instructions?
%     %
%     \answerTODO{}
%     %
%   \item Did you include the code, additional data, and instructions needed to generate
%     the figures and tables in your paper based on the raw results?
%     %
%     \answerTODO{}
%     %
%   \end{enumerate}
%   %
% \item If you used existing assets (e.g., code, data, models)\dots
%   %
%   \begin{enumerate}
%   \item Did you cite the creators of used assets?
%     %
%     \answerTODO{}
%     %
%   \item Did you discuss whether and how consent was obtained from people whose
%     data you're using/curating if the license requires it?
%     %
%     \answerTODO{}
%     %
%   \item Did you discuss whether the data you are using/curating contains
%     personally identifiable information or offensive content?
%     %
%     \answerTODO{}
%     %
%   \end{enumerate}
%   %
% \item If you created/released new assets (e.g., code, data, models)\dots
%   %
%   \begin{enumerate}
%     %
%     \item Did you mention the license of the new assets (e.g., as part of your
%     code submission)?
%     %
%     \answerTODO{}
%     %
%     \item Did you include the new assets either in the supplemental material or as
%     a \textsc{url} (to, e.g., GitHub or Hugging Face)?
%     %
%     \answerTODO{}
%     %
%   \end{enumerate}
%   %
% \item If you used crowdsourcing or conducted research with human subjects\dots
%   %
%   \begin{enumerate}
%   \item Did you include the full text of instructions given to participants and
%     screenshots, if applicable?
%     %
%     \answerTODO{}
%     %
%   \item Did you describe any potential participant risks, with links to
%     institutional review board (\textsc{irb}) approvals, if applicable?
%     %
%     \answerTODO{}
%     %
%   \item Did you include the estimated hourly wage paid to participants and the
%     total amount spent on participant compensation?
%     %
%     \answerTODO{}
%     %
%   \end{enumerate}
% \item If you included theoretical results\dots
%   %
%   \begin{enumerate}
%   \item Did you state the full set of assumptions of all theoretical results?
%     %
%     \answerTODO{}
%     %
%   \item Did you include complete proofs of all theoretical results?
%     %
%     \answerTODO{}
%     %
%   \end{enumerate}
%   %
% \end{enumerate}

\end{document}
